\documentclass[11pt,a4paper]{article}

\usepackage[utf8]{inputenc}
\usepackage[T1]{fontenc}
\usepackage[margin=1in]{geometry}
\usepackage{hyperref}
\usepackage{listings}
\usepackage{xcolor}
\usepackage{booktabs}

% ---------- Style du code ----------
\definecolor{codegray}{rgb}{0.95,0.95,0.95}
\definecolor{codeblue}{rgb}{0.13,0.29,0.53}
\definecolor{codegreen}{rgb}{0.0,0.5,0.0}
\lstset{
  backgroundcolor=\color{codegray},
  basicstyle=\ttfamily\small,
  breaklines=true,
  keywordstyle=\color{codeblue}\bfseries,
  commentstyle=\color{codegreen},
  frame=single,
  showstringspaces=false
}

% ---------- Page de titre : A COMPLETER ----------
\title{\textbf{ULK POLYTECHNIC INSTITUTE}\\
       \Large Cryptography \& Network Security (ETTCS801)\\
       \large Integrated Situation -- Final Technical Report}
\author{\textbf{Student name:} Romy \\
        \textbf{Student ID:} \dots\dots\dots\dots \\
        \textbf{Module Code:} ETTCS801 \\
        \textbf{Lecturer:} Isaac TUMWINE}
\date{\today}

\begin{document}
\maketitle
\hrule
\vspace{0.5cm}

% =====================================================
\section{Introduction \& Project Overview}
This report presents the security assessment, implementation, and traffic
filtering configuration carried out for the institute's infrastructure.
To address weak passwords, unencrypted data transfers, outdated software
and unauthorized network access, a three-tier defense strategy was applied:

\begin{enumerate}
  \item A formal risk assessment prioritizing critical assets and vulnerabilities.
  \item A Python security toolkit implementing symmetric encryption (AES via
        Fernet) and SHA-256 hashing for integrity verification.
  \item Host-based packet filtering with \texttt{iptables} on Linux to restrict
        server access according to network subnets.
\end{enumerate}

% =====================================================
\section{Risk Assessment}

\subsection{Asset and Vulnerability Identification}
\begin{itemize}
  \item \textbf{Asset 1: Student records (database/files)}
  \begin{itemize}
    \item \textit{Vulnerability:} unencrypted storage and unencrypted transfers between campuses.
    \item \textit{Consequence:} interception, tampering or exposure of grades and personal data.
  \end{itemize}
  \item \textbf{Asset 2: Central server}
  \begin{itemize}
    \item \textit{Vulnerability:} Guest network can reach the server without firewall isolation.
    \item \textit{Consequence:} network discovery, lateral movement or DoS from untrusted networks.
  \end{itemize}
  \item \textbf{Asset 3: Administrator / staff accounts}
  \begin{itemize}
    \item \textit{Vulnerability:} weak passwords and outdated software.
    \item \textit{Consequence:} brute-forcing, administrative compromise, privilege escalation.
  \end{itemize}
\end{itemize}

\subsection{Risk Evaluation and Ranking}
\begin{table}[h!]
  \centering
  \begin{tabular}{@{}llll@{}}
    \toprule
    \textbf{Risk} & \textbf{Likelihood} & \textbf{Impact} & \textbf{Priority} \\
    \midrule
    Unencrypted file storage \& transfer      & High   & High   & Critical (Rank 1) \\
    Guest network access to central server    & High   & Medium & High (Rank 2) \\
    Weak passwords \& outdated software       & Medium & High   & Medium (Rank 3) \\
    \bottomrule
  \end{tabular}
  \caption{Risk ranking matrix}
\end{table}

\begin{itemize}
  \item \textbf{Rank 1:} student records travel unencrypted over shared links;
        interception causes loss of confidentiality and reputational damage.
  \item \textbf{Rank 2:} Guest devices can scan and probe open server ports
        because no boundary firewall separates them.
  \item \textbf{Rank 3:} the impact is severe, but exploitation requires active
        reconnaissance or brute force, so it is less immediate than open
        routes and cleartext transfers.
\end{itemize}

\subsection{Recommended Security Controls}
\begin{enumerate}
  \item \textbf{Data security:} symmetric encryption (AES) for stored records
        and HTTPS/SFTP for all cross-campus transfers.
  \item \textbf{Network filtering:} \texttt{iptables} rules dropping traffic
        from the Guest subnet (\texttt{192.168.20.0/24}).
  \item \textbf{Access control:} strong password policy, multi-factor
        authentication where possible, and regular patching.
\end{enumerate}

% =====================================================
\section{Encryption and Integrity Application}

\subsection{Implementation Overview}
The utility \texttt{security\_toolkit.py} offers command-line encryption,
decryption, hashing and integrity verification using the \texttt{cryptography}
library and the Python standard library:

\begin{itemize}
  \item \textbf{Symmetric encryption:} Fernet (AES-128 in CBC mode with HMAC authentication).
  \item \textbf{Integrity hashing:} \texttt{hashlib.sha256()} produces a 256-bit
        digest before and after transfer.
  \item \textbf{Error handling:} missing files and invalid or corrupted
        ciphertexts are caught to avoid crashes.
\end{itemize}

\subsection{Usage}
\begin{lstlisting}[language=bash]
python3 security_toolkit.py --encrypt  FILE
python3 security_toolkit.py --decrypt  FILE.enc
python3 security_toolkit.py --hash     FILE
python3 security_toolkit.py --verify   FILE EXPECTED_SHA256
\end{lstlisting}

\subsection{Source Code (main functions)}
\begin{lstlisting}[language=Python]
import hashlib, os, sys
from cryptography.fernet import Fernet

KEY_FILE = "secret.key"

def generate_or_load_key():
    """Generate a Fernet key or load the existing one."""
    if not os.path.exists(KEY_FILE):
        key = Fernet.generate_key()
        with open(KEY_FILE, "wb") as kf:
            kf.write(key)
    else:
        with open(KEY_FILE, "rb") as kf:
            key = kf.read()
    return key

def calculate_sha256(file_path):
    sha256_hash = hashlib.sha256()
    with open(file_path, "rb") as f:
        for block in iter(lambda: f.read(4096), b""):
            sha256_hash.update(block)
    return sha256_hash.hexdigest()

def encrypt_file(file_path):
    fernet = Fernet(generate_or_load_key())
    with open(file_path, "rb") as f:
        data = f.read()
    with open(file_path + ".enc", "wb") as f:
        f.write(fernet.encrypt(data))

def decrypt_file(file_path):
    fernet = Fernet(generate_or_load_key())
    with open(file_path, "rb") as f:
        encrypted = f.read()
    try:
        return fernet.decrypt(encrypted)
    except Exception as e:
        print(f"[!] Decryption failed: {e}")
        sys.exit(1)
\end{lstlisting}

\subsection{Test Evidence: Encryption, Decryption and Tamper Detection}
Paste below the \emph{real} output of your terminal after running the
commands above on a sample file.

\begin{lstlisting}
$ python3 security_toolkit.py --encrypt sample_student_record.txt
<< paste your real output here >>

$ python3 security_toolkit.py --decrypt sample_student_record.txt.enc
<< paste your real output here >>

$ python3 security_toolkit.py --verify tampered_student_record.txt <hash>
<< paste your real output here >>
\end{lstlisting}

% =====================================================
\section{Network Traffic Filtering Configuration \& Testing}

\subsection{Firewall Rule Architecture}
The central server (\texttt{192.168.10.100}) uses a default-deny ingress
policy: Staff traffic (\texttt{192.168.10.0/24}) is allowed on port 80 only,
Guest traffic (\texttt{192.168.20.0/24}) is dropped, and everything else is blocked.

\begin{lstlisting}[language=bash]
sudo iptables -F
sudo iptables -P INPUT DROP
sudo iptables -P FORWARD DROP
sudo iptables -A INPUT -i lo -j ACCEPT
sudo iptables -A INPUT -m state --state ESTABLISHED,RELATED -j ACCEPT
sudo iptables -A INPUT -s 192.168.20.0/24 -j DROP
sudo iptables -A INPUT -s 192.168.10.0/24 -p tcp --dport 80 -j ACCEPT
sudo iptables -A INPUT -p tcp --dport 80 -j DROP
\end{lstlisting}

\subsection{Test Execution}
Tests were run with \texttt{netcat} (one permitted connection, two blocked).

\begin{table}[h!]
  \centering
  \small
  \begin{tabular}{@{}p{3.2cm}p{5.6cm}p{5.2cm}@{}}
    \toprule
    \textbf{Test} & \textbf{Command} & \textbf{Result} \\
    \midrule
    1. Staff, port 80 (permitted) &
      \texttt{nc -zv -s 192.168.10.15 192.168.10.100 80} &
      Connection succeeded -- \textbf{PASS} \\
    2. Guest subnet (blocked) &
      \texttt{nc -zv -w 3 -s 192.168.20.45 192.168.10.100 80} &
      Timed out (silent drop) -- \textbf{PASS} \\
    3. Staff, port 22 (blocked) &
      \texttt{nc -zv -w 3 -s 192.168.10.15 192.168.10.100 22} &
      Timed out (default DROP) -- \textbf{PASS} \\
    \bottomrule
  \end{tabular}
  \caption{Firewall test results}
\end{table}

% =====================================================
\section{Conclusion}
The measures address the vulnerabilities found in the risk assessment. Data at
rest and in transit are protected by symmetric encryption and SHA-256 integrity
checks, while the \texttt{iptables} policy isolates the central server from the
Guest subnet and from unauthorized service ports.

\vspace{0.5cm}
\noindent\textbf{GitHub repository:} \\
\url{https://github.com/Romuald-Mw/named-cryptography-network-security-exam.}

\end{document}
